\section{A uniform computable bound for weighted templates}
\label{sec:counting}

The letter sizes in a weighted template can be arbitrarily large.
We will bound its number of lattice elements using only the number of
units.  By \Cref{t:weighted-labels}, it suffices to count permutation
overgroups of the prescribed actions on each residual union of units.
Such an overgroup may have several orbits.  On each orbit we need both a
bound for the number of transitive overgroups and a bound for their
quotients by the normal closures of the prescribed actions.  The latter
bound controls how the different orbit actions can be coupled in one
subgroup.

We begin with products of natural alternating actions, allowing one
factor to act on several copies of its alphabet.  Compression in bounded
coordinate groups will let us lift a count from a block action to the
full permutation action.  We then include the private square actions
present in the weighted templates and assemble the residual orbit
actions.  All bounds are computable functions of the stated parameters;
none depends on the numerical weights of the units.

\begin{definition}[Controls]\label{n:controls}
A \emph{rank-one control} on a finite set $I$ is a group $X$ that is a
product of independent alternating groups $A_{n_i}$, with $n_i\ge7$.
Each factor acts naturally on one or more disjoint copies of its alphabet,
using the same permutation on all its copies, and fixes all other letters.
Additional fixed letters are allowed, as is the trivial group.  Its
\emph{complexity} is the number of its orbits on $I$.

A \emph{control with private squares} is the disjoint combination of a
rank-one control and square orbits, each with its own two independent
natural alternating factors.  The factors of a square act on no other
orbit.  The two alphabets of each square have the same size, at least
seven.  Its complexity again counts all orbits, including fixed letters.

A \emph{trimmed control} fixes a specified bounded number of entries in
each alphabet and retains the alternating action on the remaining entries.
If fewer than seven entries remain, the whole remaining alphabet is
declared fixed.
\end{definition}

Fixing at most $k$ entries per alphabet changes complexity by a computable
function of $k$ and the original complexity.  Indeed, each natural copy
splits into the fixed entries and at most one remaining natural orbit;
an alphabet discarded at this step has size at most $k+6$.
A private square splits into products of these alphabet orbits, so it
contributes at most $(k+6)^2$ orbits after trimming.
There is also a computable function $q_r(s)$ bounding the number of
orbits on $I^r$ for either kind of control of complexity at most $s$.
To see this, specify the orbit containing each tuple entry and then the
equality patterns of the alphabet entries in each factor.  There are at
most $2s$ factors and at most $2r$ alphabet entries to consider.
Passing from the corresponding symmetric groups to alternating groups
splits an orbit into at most $2^{2s}$ orbits.  For example, with $B_j$ the
$j$th Bell number, one may use
\begin{equation}\label{n:tuple-bound}
 q_r(s)=2^{2s}(s+1)^r B_{2r}^{\,2s}.
\end{equation}
Consequently the number of $X$-invariant $r$-ary relations is at most
$2^{q_r(s)}$.  These estimates apply to equivariant images of invariant
subsets as well: relations on such images pull back to relations on the
original tuple sets.

\begin{lemma}[Compression in a bounded base group]\label{n:compression}
Let $X$ be a rank-one control of complexity at most $s$ on $I$, and let
$D$ be a finite group of order at most $a$.  There are computable bounds,
depending only on $(s,a)$, for the following quantities:
\begin{enumerate}[label=\textup{(\roman*)}]
\item the number of $X$-invariant subgroups $J\le D^I$;
\item the number of $X$-invariant right or left cosets of each such $J$;
\item the number of tuples needed to generate $J$ under the action of
      $X$, each of which can be chosen fixed by a trimmed control of
      bounded complexity.
\end{enumerate}
Every invariant coset has a representative fixed by such a trimmed
control.  No commutativity assumption on $D$ is required.
\end{lemma}

\begin{proof}
For one factor $A_n$, combine the coordinates in its repeated natural
copies into an entry in a product $E=D^c$, where $c\le s$ is the number
of copies.  Thus the coordinates on this alphabet form a list
$(x_1,\ldots,x_n)$ in $E^n$.
Suppose first that $x$ belongs to an invariant left coset $xJ$; the case
$x\in J$ is included.  If $\sigma$ is a three-cycle of alphabet entries,
then $x^{-1}x^\sigma\in J$.  This tuple is supported on the three moved
entries, and an appropriate choice of $\sigma$ gives any prescribed
difference $x_i^{-1}x_j$ at one of them.

Let $z\in J$ be such a tuple, and let its value at a chosen entry be
$t\in E$.  Choose two fresh entries outside its support and conjugate
$z$ by a three-cycle that moves the chosen entry through those two
positions, fixing all other support entries.  If $z'$ is the resulting
tuple, then $z(z')^{-1}$ is supported on a pair and has values $t,t^{-1}$
there.  The other entries cancel coordinatewise.  Since $n\ge7$, the
fresh entries exist; since $A_n$ is two-transitive, such a pair can be
moved to any ordered pair of distinct positions.

Keep position $2$ as a reference.  For each $i\ne1,2$, use a pair tuple
in $J$ with value $x_i^{-1}x_2$ at $i$ and its inverse at position $1$.
Right multiplication changes entry $i$ to $x_2$ and accumulates the
correction at position $1$.  These are valid coordinatewise products
even when $E$ is nonabelian.  After all corrections, every entry except
possibly the first is the common reference value.  Do this independently
on every alphabet.  The resulting representative is constant beyond
the first two positions of every alphabet and is fixed by the control
fixing those positions.

Let $\mathcal C$ be the set of all tuples in $D^I$ constant beyond the
first two entries on each alphabet.  The fixed coordinates and the three
values per natural copy give, for example, $|\mathcal C|\le a^{3s}$.
Every pair multiplier used above is an $X$-translate of an element of
$J\cap\mathcal C$.  Therefore $J$ is generated by the $X$-translates of
$J\cap\mathcal C$.  This bounds its generating list and the number of
possible $J$'s, for instance by $2^{|\mathcal C|}$.  Every invariant
left coset meets $\mathcal C$, so there are at most $|\mathcal C|$
such cosets for fixed $J$.  Inversion gives the corresponding assertion
for right cosets.  The required trimming bounds were established after
\Cref{n:controls}.
\end{proof}

\begin{lemma}[Lifting across bounded fibers]\label{n:lifting}
Let $X$ be a rank-one control of complexity at most $s$ on $I$, let $D$
have bounded order, and let $U'\le S_I$ contain $X$.  Suppose $U'$ is
generated over $X$ by a bounded
list of elements $g$ such that each $X\cap gXg^{-1}$ contains a rank-one
control of bounded complexity.  Then the number of subgroups
\[
 V\le D^I\rtimes U',\qquad \operatorname{proj}(V)=U',
 \qquad (1,X)\le V,
\]
is bounded computably in $s$ and these parameters.  If
$[U':\langle X^{U'}\rangle]$ is bounded, then so is
$[V:\langle(1,X)^V\rangle]$.

In the permutation action with bounded fibers, each such $V$ is generated
over its pure copy of $X$ by boundedly many elements whose intersections
with that pure control contain rank-one controls of bounded complexity
on the expanded point set.
\end{lemma}

\begin{proof}
The base kernel $J=V\cap D^I$ is $X$-invariant, so its possibilities are
bounded by \Cref{n:compression}.  For a prescribed generator $g$, write
the fiber of the projection as $(Jf,g)$.  If
$h\in X\cap gXg^{-1}$, multiply a lift on the left by $(1,h)$ and on
the right by the matching pure lift of $g^{-1}h^{-1}g$.  The resulting
element still lies above $g$, and its base entry is $h\cdot f$.
Thus $Jf$ is invariant under the specified intersection control.
\Cref{n:compression}, applied to that control, bounds its possible
cosets and supplies a representative $f$ fixed by a trimmed subcontrol.
The kernel, these lifts, and $(1,X)$ generate $V$, proving the count.

Put $N=\langle(1,X)^V\rangle$.  In $V/N$, all $X$-translates of a
kernel generator have the same image.  Hence the image of $J$ is already
generated by the images of a bounded list of tuples, with no translates
required.  A subgroup of $D^I$ generated by $k$ tuples has order at most
\begin{equation}\label{n:base-order}
 |D|^{\,|D|^k}.
\end{equation}
Indeed, each coordinate has a type in $D^k$, and coordinates of the same
type evaluate every word in the generators identically.  Projection to
one coordinate of each type embeds the subgroup into at most $|D|^k$
copies of $D$.  The quotient of $V/N$ by the image of $J$ is
$U'/\langle X^{U'}\rangle$, which proves the index bound.

For the last assertion, let $h$ belong to the trimmed intersection
control fixing the chosen $f$.  Conjugating its pure lift by $(f,g)$
takes it back into pure $X$: the base term disappears because $h$ fixes
$f$, and $g^{-1}hg\in X$.  The compressed kernel generators likewise
commute with their trimmed pure controls.  A natural action repeated
over bounded fibers is still a rank-one control, with complexity
multiplied by a bounded factor.  These observations give all the
asserted intersections on the expanded set and remain valid after a
bounded number of such lifts.
\end{proof}

We use the following consequence of the primitive order theorem with an
explicit effective threshold.  Mar\'oti's Theorem
\citep[Theorem~1.1]{Maroti2002} says that a primitive group of degree $d$
is either one of the four Mathieu exceptions, has order at most
$d^{1+\log_2d}$, or lies between
\[
 A_\ell^e\quad\hbox{and}\quad S_\ell\wr S_e
 \quad\hbox{on degree}\quad d=\binom{\ell}{h}^{e}.
\]
Replacing subsets by their complements permits $h\le\ell/2$.

\begin{lemma}[Large primitive parts]\label{n:primitive-parts}
For each positive integer $s$ there is a computable threshold $H(s)$
with the following properties for $d>H(s)$.  A primitive group of degree
$d$ and order at least $\lfloor\sqrt{d/s}\rfloor!/2$ has one of the
alternating-product actions
\begin{equation}\label{n:three-actions}
 (e,h)=(1,1),\qquad(1,2),\qquad(2,1),
\end{equation}
with $\ell$ arbitrarily large as the threshold is increased.  Its socle
is $A_\ell^e$ and has index at most eight.  If its order is at least
$\lfloor d/s\rfloor!/2$, it contains $A_d$ in the natural action.
One may take
\begin{equation}\label{n:threshold}
 H(s)=2^{\,2^{(s+100)^{10}}}.
\end{equation}
\end{lemma}

\begin{proof}
Write $p=\log_2d$.  Outside the bounded Mathieu exceptions, the small-order
alternative has logarithm at most $p(1+p)$.  In an alternating-product
action, $e\le p$.  Unless one of \Cref{n:three-actions} holds,
$\ell\le4d^{1/3}$: if $e\ge3$, use $\binom\ell h\ge\ell$; if
$e=2,h\ge2$, use $\binom\ell h\ge\binom\ell2$; and if $e=1,h\ge3$,
use $\binom\ell h\ge\binom\ell3$.  Consequently the logarithm of
$(\ell!)^e e!$ is at most $10p^2d^{1/3}$.  For the two additional
possibilities $(1,2)$ and $(2,1)$, the same estimate with $d^{1/2}$ in
place of $d^{1/3}$ is an upper bound.

Under \Cref{n:threshold}, we have $s<p$ and
$2^{p/6}>100p^4$.  Compare the preceding upper bounds with
\[
 \log_2(v!/2)\ge v-2,
 \qquad
 v=\lfloor\sqrt{d/s}\rfloor
 \quad\hbox{or}\quad v=\lfloor d/s\rfloor,
\]
respectively.  The square-root factorial excludes the small-order
alternative and all cases outside \Cref{n:three-actions}; the linear
factorial excludes the other two actions as well.  The socles in the
three remaining cases are respectively $A_\ell$, $A_\ell$, and
$A_\ell^2$.  Their full indicated overgroups have quotient orders two,
two, and eight.  In the last case primitivity forces transitivity on
the two slots: fixing the slots would preserve the partition into
lines in either direction.  The centralizer of a natural alternating
group, and of its action on two-subsets, is trivial; a commuting
permutation is determined by the image of one point, whose stabilizer
fixes no other point in these actions.  In the product action,
commuting with each independently supported factor likewise fixes
each coordinate.  Finally, the intersection of a normal subgroup with the
alternating product is a subproduct of its simple factors, and slot
transitivity forces every nontrivial such intersection to be
the whole product.  A normal subgroup with trivial intersection
centralizes that product and is therefore trivial.  This proves the
assertions about the socle and unique minimal normal subgroup.
\end{proof}

The next overgroup argument needs bijections between equal-sized parts
that commute with one common subcontrol.  For a rank-one control fixing
each part setwise, its \emph{profile} on a part records the number of
fixed letters and the number of natural copies of each factor of that
control.  Equal profiles give equivariant bijections by matching copies
of the same factor with their alphabet coordinates, and matching the
fixed letters.  Equal cardinalities alone need not give equal profiles;
the following lemma supplies a suitable subcontrol.

\begin{lemma}[Effective equal profiles]\label{n:profiles}
Suppose a finite set is partitioned into at most $s$ equal-sized parts
$B$, each preserved by a rank-one control $X$.  Write the size of $B$ as
\[
 b_B+\sum_i k_{Bi}n_i,
\]
where $b_B$ counts its fixed letters, $k_{Bi}$ counts its natural copies
of the factor $A_{n_i}$ of $X$, and the number of factors and all
nonnegative integers $b_B,k_{Bi}$ are at most $s$.  There is a rank-one
subcontrol of $X$, of complexity bounded computably in $s$, that has
identical profiles on all parts.  If the common part size is at least
two, the subcontrol may also be chosen with at least two fixed letters
in every part.
\end{lemma}

\begin{proof}
The empty set is immediate, so suppose there is at least one part.
Subtract the expression for one part from each other expression.  This
gives a system $An=c$ with bounded integer coefficients and bounded
numbers of equations and variables.  No equation prescribes the
common total size.  Every nonnegative integer solution admits a
decomposition
\begin{equation}\label{n:profile-decomposition}
 n=u+\sum_{w\in\mathcal H}m_w w,\qquad m_w\in\mathbb Z_{\ge0},
\end{equation}
where $u$ belongs to the finite set of coordinatewise minimal
nonnegative solutions of $Au=c$, and $\mathcal H$ is the finite set
of minimal nonzero nonnegative solutions of $Aw=0$.
The finiteness is Dickson's Lemma \citep{Dickson1913}.  Given $n$, choose
such a $u\le n$; then $n-u$ is homogeneous.  Unless it is zero, subtract
a minimal nonzero homogeneous solution lying below it and repeat.
The sum of its coordinates strictly decreases, proving the decomposition.

These finite lists can be found effectively.  Enumerate solutions in
larger and larger boxes and test whether their upward orthants cover
all nonnegative solutions, using the first-order theory of addition
and order \citep{Presburger1929}.  The test has only fixed integer
coefficients: for a proposed particular list it asks whether
\[
 \forall v\in\mathbb Z_{\ge0}^r\quad
 (Av=c\ \Longrightarrow\ \bigvee_u u\le v),
\]
and the homogeneous test adds $v\ne0$.  Dickson's Lemma guarantees
termination; discarding nonminimal listed entries gives the desired
lists.  There are only boundedly many possible coefficient systems,
so the lengths and coordinates of all these lists have computable
uniform bounds.

For each original alphabet $i$, fix $u_i$ entries and split the rest
into $w_i$ sets of size $m_w$ for each $w\in\mathcal H$.
For each $w$, use one common natural alternating factor $A_{m_w}$
diagonally on all these sets, across all original alphabets and repeated
copies.  Thus the profiles on different parts refer to the same actual
factors, even when two factors have equal degrees.  If $m_w<7$, fix
those entries instead.  The number of copies on part $B$ for the new factor
is $\sum_i k_{Bi}w_i$, independent of $B$ because $Aw=0$.
The remaining fixed-letter count is also independent of $B$, by
$Au=c$.  This proves identical profiles and the complexity bound.
Suppose the common part size is at least two.  The number of fixed
letters is already the same in every part.  If it is less than two,
some new factor $A_m$, with $m\ge7$, acts on a part.  By the identical
profiles it has the same positive number of natural copies on every
part.  Fix two entries of its alphabet, simultaneously on all these
copies.  If $m\ge9$, retain $A_{m-2}$ on the remaining entries.  If
$m=7$ or $8$, declare all entries of these copies fixed, as required
by the trimming convention.  This adds at most eight fixed letters
per copy, and the number of copies is computably bounded by the
preceding construction.  The new profiles are still identical, and
every part now has at least two fixed letters.  If no factor acts,
all letters were already fixed, so the same conclusion follows from
the common part size.  The trivial control is allowed; discarding
the last active factor therefore causes no exception.  This proves
the fixed-letter assertion with computably bounded complexity.
\end{proof}

\begin{proposition}[Rank-one overgroups]\label{n:rank-one}
There are computable functions of $s$ that bound, for every rank-one
control $X$ of complexity at most $s$ on a finite nonempty set $I$,
\begin{enumerate}[label=\textup{(\roman*)}]
\item the number of transitive groups $U$ with $X\le U\le S_I$ and
      the index $[U:\langle X^U\rangle]$;
\item the length of a generating list for $U$ over $X$, chosen so
      that every generator $g$ has a rank-one control of computably
      bounded complexity inside $X\cap gXg^{-1}$.
\end{enumerate}
\end{proposition}

\begin{proof}
We prove a more precise assertion with a second parameter $h$ bounding
the number of strict steps in any chain of $U$-invariant equivalences.
Every such equivalence is a union of $X$-orbits on $I^2$, so initially
$h\le q_2(s)$ suffices.  We induct on $h$, allowing the complexity to
change by computable functions at a step.  If $|I|=1$, the assertion
is immediate.  Otherwise choose a minimal nondiscrete $U$-invariant
partition $\mathcal B$; the universal partition is allowed.  Its
parts have a common size $d\ge2$, and a part stabilizer acts
primitively on its part.  There are at most $2^{q_2(s)}$ possible
partitions.  Let $I'$ be the set of parts.

\emph{A natural copy maps injectively to $I'$.}
The map from an $X$-natural copy to $I'$ is either constant or
injective, since a natural alternating action is primitive.
If a copy for one factor maps injectively, all letters mapping to
that moving orbit of parts come from injective copies of the same
factor.  Indeed a copy on which this factor is trivial cannot have
an equivariant image in its nontrivial natural orbit.  Each
contributing copy supplies only one letter to a part, whence $d\le s$.

Retain only the natural factors of degree greater than $d+2$ to form
a smaller control $X_0$.  Discarded factors have bounded degree,
so $X_0$ has bounded complexity.  It induces a rank-one control
$X_0'$ on $I'$, whose complexity is no larger than that of $X_0$.
A part stabilizer in $X_0$ is a product of alternating
groups or natural point stabilizers.  Each nontrivial factor is
simple and has order greater than $d!$, so its action on a fiber of
size $d$ is trivial.  Label one fiber in each $X_0$-orbit of parts
and transport the labeling by $X_0$.  The trivial stabilizer action
makes this well-defined.  This labels $X_0$ as a pure control in a
bounded-fiber wreath product, and the labeling can be fixed from
$X_0$ and $\mathcal B$, independently of $U$.

The projected transitive group $U'$ has equivalence-chain bound at
most $h-1$: lift a chain on $I'$ and append the discrete partition
of $I$ at its fine end.  Apply the induction hypothesis to
$(X_0',U')$, then \Cref{n:lifting} with fiber group $S_d$.
This bounds the choices of $U$, its index modulo the normal closure
of $X_0$, and its generators with the required intersection controls
on $I$.  They also prove the assertions for $X\ge X_0$: its normal
closure is larger, generation over it needs no additional elements,
and the intersection controls remain contained in
$X\cap gXg^{-1}$.

\emph{Every natural copy maps constantly to $I'$.}
Now $X$ fixes every part, so their number $b$ is at most $s$.
If $d\le H(s)$, then $|I|\le sH(s)$ is computably bounded, and
finite subset counts bound every subgroup choice and generator
list.  The trivial subgroup is an admissible intersection control
of bounded complexity on this bounded domain.  Thus assume $d>H(s)$.
Each part contains at most $s$ natural copies or fixed letters,
so some natural axis of degree at least $d/s$ occurs in it.
\Cref{n:primitive-parts} implies that its primitive part-stabilizer
action contains $A_d$.

Let $K$ be the kernel of the action on $I'$ and put
\[
 T_0=K\cap\prod_{B\in\mathcal B}A_B.
\]
The projection of $K$ to each part is normal in its primitive
part-stabilizer group and contains the nontrivial axis action.
It therefore contains $A_B$.  The commutator subgroup of $K$ lies
in $T_0$ and still projects onto every $A_B$.  Thus $T_0$ is a
subdirect product of these nonabelian simple groups, hence a
product of full twisted diagonals over a partition of the parts.
It is normal in $U$.

For completeness, such a subdirect product is self-normalizing
inside the direct product of its simple factors.  On each diagonal
class, an element normalizing the diagonal has coordinate
conjugations that agree through all the tying isomorphisms.  Its
coordinates therefore differ from a member of that diagonal by
elements of the centers, which are trivial.  Consequently the
whole block-stabilizer normalizer satisfies
\begin{equation}\label{n:rank-one-normalizer}
 \bigl[N_{\,S_d\wr S_b}(T_0):T_0\bigr]\le2^b b!.
\end{equation}
In high degree, automorphisms of $A_d$ are induced uniquely by
permutations of its natural alphabet.  Record each twisted tie
by matching the corresponding letters across its tied parts.
The resulting relation on $I^2$ is $X$-invariant, since $U$
normalizes $T_0$.  It determines the ties and their twists, so
there are boundedly many choices of $T_0$.  For fixed $T_0$,
\Cref{n:rank-one-normalizer} bounds all choices of $U$ by subset
counts in its quotient.

Here $1<X\le T_0$: $X$ fixes each part and its natural alternating
actions are even on it.  Since $U$ is transitive on parts, it is
transitive on the diagonal simple factors of $T_0$.  A normal
subgroup of their product is a subproduct.  Thus
$\langle X^U\rangle$ contains all of $T_0$, giving the index
bound in (i).

To obtain (ii), first generate $T_0$ by supported three-cycles in
its diagonal factors.  Such a cycle is specified by at most $3b$
letters, recording its simultaneous support in all tied parts.
There are boundedly many $X$-orbits of these specifications by
\Cref{n:tuple-bound}.  Take one representative from each orbit.
Each commutes with the control trimmed to fix its support, and
these representatives generate $T_0$ over $X$.

It remains to choose suitable representatives for $U/T_0$.
Write the size of each part as
$b_B+\sum_i k_{Bi}n_i$, with fixed-letter counts and copy
multiplicities for $X$.  Their coefficients are bounded by $s$.
\Cref{n:profiles} gives a subcontrol $X_1\le X$ of bounded
complexity, with identical profiles and at least two fixed
letters on every part.  Fix a coset in $U/T_0$.  By multiplying
its representative by $T_0$, the map from one representative
part in each diagonal class to its destination may be changed
to any bijection of the prescribed parity.  The maps on its
other tied parts then correspond through the tying
identifications, since the representative transports the
diagonals.  Those identifications are $X$-equivariant, because
$X\le T_0$.  The identical $X_1$-profiles allow the chosen
bijections to be $X_1$-equivariant; a swap of two fixed letters
corrects their parity if necessary.  The resulting coset
representative centralizes $X_1$, so
$X_1\le X\cap gXg^{-1}$.  There are boundedly many cosets by
\Cref{n:rank-one-normalizer}.

This completes the induction.  All coefficient lists used for
profiles are drawn from bounded ranges, and all bounded-domain
branches use finite subset counts.  The only recursive call
strictly decreases $h$, even when the complexity grows after
small axes are discarded.  Hence the bounds proved are computable
functions of $(s,h)$, and substituting $h=q_2(s)$ proves the
Proposition.
\end{proof}

\begin{proposition}[Overgroups of controls with private squares]
\label{n:private-squares}
There are computable functions $B(s),C(s)$ such that every control
$X$ with private squares of complexity at most $s$ on a finite
nonempty set $I$ has at most $B(s)$ transitive overgroups in $S_I$,
and every such overgroup $U$ satisfies
\[
 [U:\langle X^U\rangle]\le C(s).
\]
The functions can be taken nondecreasing and at least one.
\end{proposition}

\begin{proof}
We count three successive choices: the block partition, the local
primitive structures and their diagonal ties, and the overgroup in the
resulting wreath product with bounded local groups.  The first two
choices will be encoded by
invariant relations on tuples of bounded length.

Again choose a minimal nondiscrete $U$-block partition $\mathcal B$,
with part size $d\ge2$ and part set $I'$.  The case $|I|=1$ is
immediate, and the choices of $\mathcal B$ are bounded by
\Cref{n:tuple-bound}.

The restriction of this partition to a square orbit is either
constant, a partition into lines in one direction, or discrete.
Indeed the point stabilizer in the product of its two alternating
factors is the product of two maximal nonnormal natural point
stabilizers.  If an overgroup of that stabilizer projects larger
in an axis, conjugating the independently present point
stabilizer by a mover generates that entire alternating factor
independently.  These are exactly the four resulting block
possibilities.  The discrete possibility cannot occur here:
privacy means that no other $X$-orbit can map equivariantly onto
the moving orbit of parts of this square, so a part meeting it
would have size one, contrary to $d>1$.

It follows that the image $X'$ on $I'$ is rank-one.  Each moving
part orbit is identified equivariantly with the natural alphabet
of its contributing axis.  Distinct such orbits for the same
axis are diagonal natural copies; the factor is trivial on
every other part orbit, and the active factors remain independent.
There are at most $s$ part orbits.  Moreover, an injective
rank-one copy on $I'$ can only share its parts with copies of
the same factor, one letter per copy.  Other orbits, including
all the private squares, have trivial action of that factor.
Thus an injective rank-one copy still forces $d\le s$.

If $d$ is bounded in terms of $s$, every private square has side
at most $d$, since one of its axis lines lies in a part.  Discard
its factors and declare its letters fixed.  There are boundedly
many such letters, and what remains is a rank-one control of
bounded complexity on $I$.  \Cref{n:rank-one} bounds the desired
number and index, also for the original larger $X$.

Assume henceforth that $d>H(s)$, increasing the displayed
threshold if needed only within a computable bound.  In a part
$B$, each original orbit contributes a whole orbit, an axis
line, or a singleton.  One contribution has size at least
$d/s$.  Therefore some axis of degree at least
$\sqrt{d/s}$ acts nontrivially inside $B$.  It lies in the
kernel of the action on $I'$: square axes are private, and an
injective rank-one occurrence of its factor anywhere would
force $d\le s$.

\paragraph{The local socle and its normal closure.}
By \Cref{n:primitive-parts}, the socle $S_B$ of the primitive
part-stabilizer group is $A_\ell^e$ in one of
\Cref{n:three-actions}.  It is the unique minimal normal
subgroup, has trivial centralizer, and its quotient has order
at most eight.  Let $K$ be the part kernel.  Its projection to
$B$ is a nontrivial normal subgroup of the primitive group,
so contains $S_B$.  Set
\[
 T_0=K\cap\prod_{B\in\mathcal B}S_B.
\]
This intersection projects onto every $S_B$.  To verify this
without a dependence on the number of parts, take the
$8!$-th powers of all elements of $K$; these lie in $T_0$.
In an individual simple site, every $11$-cycle is an $8!$-th
power, since $11\nmid8!$.  Such cycles generate the large
alternating group.  Since the projection of $K$ contains the
whole local socle, these powers project to a generating set
of every individual site, hence onto $S_B$.

Thus $T_0$ is a subdirect product on all the simple sites,
and is a product of full twisted diagonals.  Uniqueness of
the local socles makes it normal in $U$.  The local slots are
transitively permuted by a part stabilizer, as required by
primitivity in the two-slot case; $U$ is also transitive on
parts.  Hence $U$ is transitive on all simple sites and on
the diagonal simple factors of $T_0$.  The kernel axis in
$X$ found above has a nontrivial local projection, and the
local centralizer of $S_B$ is trivial.  Commuting that axis
with $T_0$ therefore gives a nonidentity element in
$T_0\cap\langle X^U\rangle$.  This intersection is normal
in $U$, so transitivity on the diagonal factors gives
\begin{equation}\label{n:socle-normal-closure}
 T_0\le\langle X^U\rangle.
\end{equation}

\paragraph{Encoding structures and diagonal ties.}
We next count the local structures and their diagonal ties
without assuming that $|I'|$ is bounded.  In the natural
case $(1,1)$, the site alphabet is the set of part letters.
In case $(1,2)$, choose the graph in which two distinct
two-subsets are adjacent when they intersect.  In case
$(2,1)$, choose adjacency when two distinct points share a
coordinate.  Choose this structure on one part and
transport it by $U$ to all parts; the initial part
stabilizer preserves it, so transport is well-defined.
The resulting relation on $I$ is $X$-invariant and hence
has boundedly many possibilities.

These graphs recover their natural alphabets.  In the two-subset graph,
more than three pairwise-intersecting distinct pairs have a common entry;
the large maximal cliques are exactly the stars, one for each alphabet
entry.  An ordered adjacent pair determines the star at its common entry.
In the sharing graph, a clique of size at least two lies on a unique
coordinate line, so an edge determines that line.  Disjointness groups
the lines into two directions; the lines in each direction form one
recovered site alphabet.  Thus the recovered alphabet points are stars
in the two-subset case and coordinate lines in the product case.
Their disjoint union is an equivariant image of the $X$-invariant subset
of $I^2$ consisting of ordered edges within the parts.  In the natural
case use the diagonal pairs $(p,p)$ instead.  Relations on pairs of
recovered alphabet points therefore pull back to relations on $I^4$,
whose orbit count is bounded by \Cref{n:tuple-bound}.

An isomorphism between these part structures is induced
by natural alphabet bijections, with a possible
permutation of the two slots.  Isomorphisms between high
alternating simple factors are likewise induced uniquely
by alphabet bijections.  Thus record the twisted
diagonals of $T_0$ by their matched-letter relation on
the recovered alphabets.  This relation is equivariant,
because $X$ normalizes $T_0$, and it determines the full
twisted diagonals.  Its pullback is an invariant relation
on $I^4$, so \Cref{n:tuple-bound} bounds all choices of
$T_0$ for the fixed block structure.  This explicit
four-tuple bound also covers ties between sites in
different parts.

\paragraph{Bounded-group lifts.}
Let $W$ be the stabilizer of the partition and chosen
part structures.  In fixed local coordinates its
homomorphism
\begin{equation}\label{n:normalizer-wreath}
 W\longrightarrow D^{I'}\rtimes S_{I'}
\end{equation}
records the alphabet parities and slot permutations,
where $|D|\le8$.  Its kernel is $\prod_B S_B$.
On $N_W(T_0)$ the kernel is exactly $T_0$, because a
subdirect product of nonabelian simple groups is
self-normalizing inside their direct product, as
proved in \Cref{n:rank-one}.  Hence $N_W(T_0)/T_0$
embeds in the bounded-group wreath product in
\Cref{n:normalizer-wreath}.

We can make the image of $X$ pure in this wreath product.
A part stabilizer in $X$ is a product of large
alternating groups, allowing one natural point
stabilizer $A_{n_i-1}$ when the part orbit is natural.
Since $n_i\ge7$, each nontrivial factor is simple of
order greater than eight.  Its local homomorphism to
$D$ is therefore trivial.  Fix a local quotient frame
at one representative of each $X$-orbit of parts and
transport it by $X$.  The stabilizer property makes
this unambiguous; changing frames amounts only to
changing coordinates by local elements of $D$.
All base components of $X$ are now identity.  These
frames depend only on the structures and $X$, not
on the subgroup $U$ subsequently being counted.

Apply \Cref{n:rank-one} to the rank-one part action
$X'$ on $I'$, of complexity at most $s$.  It bounds its possible transitive
overgroups and supplies the generators with
intersection controls needed in \Cref{n:lifting}.
With this complexity bound and $|D|\le8$, that lemma bounds the possible images
$U/T_0$ in \Cref{n:normalizer-wreath} and their indices
modulo the normal closure of the pure image of $X$.
For fixed structures and $T_0$, the image determines
$U$ uniquely, because $U$ contains the kernel $T_0$.
By \Cref{n:socle-normal-closure}, these quotient indices
are exactly $[U:\langle X^U\rangle]$.  Combining the
bounded choices of partitions, structures, matching
relations and lifts proves both assertions.
\end{proof}

\begin{theorem}[Uniform computable template bound]\label{n:template-bound}
There is a total computable nondecreasing function
$\Phi:\mathbb Z_{\ge0}\to\mathbb Z_{\ge1}$ such that every weighted
template
\[
 F=[S_Z\times Y,S_{Z\sqcup(\bigcup T)}]
\]
of \Cref{sec:weighted}, with at most $t$ units, excluding
the foundation $Z$ from this count, and
$|Z|>|\bigcup T|$, satisfies $|F|\le\Phi(t)$.
The bound is independent of all full-unit sizes, square sides,
auxiliary axis lengths, offset sizes, and the size of $Z$.
No stock-size hypothesis is required for this counting assertion.
\end{theorem}

\begin{proof}
On any nonempty union of units, the prescribed product $Y$
contains a control with private squares whose complexity is
bounded computably in the unit count.  On a full unit of size
at least seven, keep its natural alternating group; a smaller
full unit contributes at most six fixed letters.  On an
auxiliary unit $(x/4)\times2\times2$, if $x/4\ge7$, keep
$A_{x/4}$ on its four repeated natural copies.  Otherwise
its domain has at most $24$ letters and may be fixed entirely.
On each primary square keep its two prescribed alternating
factors; their privacy is part of the independent unit
construction.  Thus complexity at most $24t$ suffices.

Consider an arbitrary residual overgroup $U\ge Y_R$ on a
union $R$ of at most $t$ units.  Its orbits are unions of
the transitive units, so there are at most $B_t$ choices
for this orbit partition.  On every orbit $O$, its
transitive projection $U_O$ has at most $B(24t)$ choices,
by \Cref{n:private-squares}.  Let $N_O$ be the normal
closure of the prescribed unit groups on $O$ in $U_O$.
Their contained control shows that
$[U_O:N_O]\le C(24t)$.

The subgroup $N_O$, supported independently on $O$, is
actually contained in $U$: its initial unit groups are
independently supported, and conjugating them by $U$
realizes their normal closure in the projection $U_O$
without adding support on another orbit.  Consequently
$U$ contains $\prod_O N_O$.  Once its projections are
fixed, its remaining choice is a subgroup in
$\prod_O(U_O/N_O)$, a group of order at most
$C(24t)^t$.  There are at most $2^{C(24t)^t}$ such
subgroups.  The empty residual set has one choice.

By the exact literal labeling of the weighted template,
each lattice element consists of an absorbed subset
$I\subseteq T$ and one of these residual overgroups on
$T\setminus I$.  For $t\ge1$, therefore, one may take
the increasing envelope of
\begin{equation}\label{n:explicit-template-envelope}
 2^t B_t\,B(24t)^t\,2^{C(24t)^t},
\end{equation}
and put $\Phi(0)=1$.

To check effectivity of the whole construction, every
invariant-relation choice used above is bounded by a
power of two with exponent a tuple-orbit bound from
\Cref{n:tuple-bound}.  Every bounded quotient has at
most as many subgroups as subsets.  Trimming and small
degree branches use computable finite ranges; equal
profiles use maxima over finitely many bounded
coefficient systems and the terminating Presburger
procedure in \Cref{n:profiles}.  The recursion of
\Cref{n:rank-one} strictly decreases the equivalence
chain parameter.  These operations define total
computable bounds and never take a maximum by a
finite-group search over unbounded unknown degrees.
\end{proof}
